% GENERATED FILE. Edit canonical lessons, then run npm run generate:books.
% canonical-source-hash: fnv1a64:736a8049569ffdd7
% canonical-lessons: MR-C518-yantra, MR-C518-upakaran, MR-C518-yantramanav, MR-C518-svayancalit, MR-C518-phavarne

\chapter{Machines and Devices}
\label{ch:mr-machines-and-devices}
\hlchaptermodality{518}

\begin{chapteropening}
\textbf{By the end of this chapter:} \emph{I can talk about machines, devices and robots.}

Complete the last lesson of chapter 518: I can talk about machines, devices and robots.
\end{chapteropening}

\section[\texorpdfstring{\mr{यंत्र} (yantra)}{यंत्र (yantra)}]{\mr{यंत्र} (yantra) — a machine}

\label{lesson:MR-C518-yantra}



\begin{quote}
Before the new one: say the Marathi for a link, then the Marathi for to store.
\end{quote}

\subsection*{You'll want to know: \mr{यंत्र}}
\textbf{\mr{यंत्र}} — \emph{yantra} — \textquotedblleft{}a machine\textquotedblright{}.

\begin{quote}
\mr{हे} \mr{यंत्र} \mr{विजेवर} \mr{चालतं}.
\end{quote}

\begin{quote}
\emph{he yantra vijevar cāltaṁ.}
\end{quote}

\begin{quote}
This machine runs on electricity.
\end{quote}

Neuter.

\begin{cousinweb}
\textbf{\mr{यंत्र}} (\emph{yantra}) is a tatsama, taken from Sanskrit \emph{yantra}, \textquotedblleft{}an instrument, a machine\textquotedblright{} (Monier-Williams).
\end{cousinweb}

\begin{grammarlens}[title={What you've built}]
One more word for describing what you lived and felt.
\end{grammarlens}

\subsection*{Your turn}
\begin{itemize}
\raggedright
  \item \emph{Say it:} \emph{yantra}
  \item \emph{Say it:} \emph{yantra}, once more
  \item \emph{Say it:} \emph{he yantra vijevar cāltaṁ.} — This machine runs on electricity.
  \item \emph{Recall:} say the Marathi for a link, then the Marathi for to store, then say \emph{yantra} again
\end{itemize}

\begin{tcolorbox}[breakable,colback=teal!4,colframe=teal!35!black,title={Before you move on}]
What does \mr{यंत्र} mean? (\textquotedblleft{}A machine\textquotedblright{}.) Say it once more.
\end{tcolorbox}
\section[\texorpdfstring{\mr{उपकरण} (upakaraṇ)}{उपकरण (upakaraṇ)}]{\mr{उपकरण} (upakaraṇ) — a device; an appliance}

\label{lesson:MR-C518-upakaran}



\begin{quote}
Before the new one: say the Marathi for to store, then the Marathi for a machine.
\end{quote}

\subsection*{You'll want to know: \mr{उपकरण}}
\textbf{\mr{उपकरण}} — \emph{upakaraṇ} — \textquotedblleft{}a device; an appliance\textquotedblright{}.

\begin{quote}
\mr{स्वयंपाकघरात} \mr{विजेची} \mr{खूप} \mr{उपकरणं} \mr{आहेत}.
\end{quote}

\begin{quote}
\emph{svayampākgharāt vijecī khūp upakaraṇaṁ āhet.}
\end{quote}

\begin{quote}
There are many electric appliances in the kitchen.
\end{quote}

Neuter. Plural: \mr{उपकरणं}, \emph{upakaraṇaṁ}.

\begin{cousinweb}
\textbf{\mr{उपकरण}} (\emph{upakaraṇ}) is a tatsama, taken from Sanskrit \emph{upakaraṇa}, \textquotedblleft{}an instrument, a means\textquotedblright{} (Monier-Williams).
\end{cousinweb}

\begin{grammarlens}[title={What you've built}]
One more word for describing what you lived and felt.
\end{grammarlens}

\subsection*{Your turn}
\begin{itemize}
\raggedright
  \item \emph{Say it:} \emph{upakaraṇ}
  \item \emph{Say it:} \emph{upakaraṇ}, once more
  \item \emph{Say it:} \emph{svayampākgharāt vijecī khūp upakaraṇaṁ āhet.} — There are many electric appliances in the kitchen.
  \item \emph{Recall:} say the Marathi for to store, then the Marathi for a machine, then say \emph{upakaraṇ} again
\end{itemize}

\begin{tcolorbox}[breakable,colback=teal!4,colframe=teal!35!black,title={Before you move on}]
What does \mr{उपकरण} mean? (\textquotedblleft{}A device; an appliance\textquotedblright{}.) Say it once more.
\end{tcolorbox}
\section[\texorpdfstring{\mr{यंत्रमानव} (yantramānav)}{यंत्रमानव (yantramānav)}]{\mr{यंत्रमानव} (yantramānav) — a robot}

\label{lesson:MR-C518-yantramanav}



\begin{quote}
Before the new one: say the Marathi for a machine, then the Marathi for a device.
\end{quote}

\subsection*{You'll want to know: \mr{यंत्रमानव}}
\textbf{\mr{यंत्रमानव}} — \emph{yantramānav} — \textquotedblleft{}a robot\textquotedblright{}.

\begin{quote}
\mr{हा} \mr{यंत्रमानव} \mr{घर} \mr{झाडतो}.
\end{quote}

\begin{quote}
\emph{hā yantramānav ghar jhāḍto.}
\end{quote}

\begin{quote}
This robot sweeps the house.
\end{quote}

Masculine.

\begin{cousinweb}
\textbf{\mr{यंत्रमानव}} (\emph{yantramānav}) is a modern tatsama coinage from Sanskrit \emph{yantra}, \textquotedblleft{}machine\textquotedblright{}, and \emph{mānava}, \textquotedblleft{}human\textquotedblright{} (Monier-Williams).
\end{cousinweb}

\begin{grammarlens}[title={What you've built}]
One more word for describing what you lived and felt.
\end{grammarlens}

\subsection*{Your turn}
\begin{itemize}
\raggedright
  \item \emph{Say it:} \emph{yantramānav}
  \item \emph{Say it:} \emph{yantramānav}, once more
  \item \emph{Say it:} \emph{hā yantramānav ghar jhāḍto.} — This robot sweeps the house.
  \item \emph{Recall:} say the Marathi for a machine, then the Marathi for a device, then say \emph{yantramānav} again
\end{itemize}

\begin{tcolorbox}[breakable,colback=teal!4,colframe=teal!35!black,title={Before you move on}]
What does \mr{यंत्रमानव} mean? (\textquotedblleft{}A robot\textquotedblright{}.) Say it once more.
\end{tcolorbox}
\section[\texorpdfstring{\mr{स्वयंचलित} (svayañcalit)}{स्वयंचलित (svayañcalit)}]{\mr{स्वयंचलित} (svayañcalit) — automatic}

\label{lesson:MR-C518-svayancalit}



\begin{quote}
Before the new one: say the Marathi for a device, then the Marathi for a robot.
\end{quote}

\subsection*{You'll want to know: \mr{स्वयंचलित}}
\textbf{\mr{स्वयंचलित}} — \emph{svayañcalit} — \textquotedblleft{}automatic\textquotedblright{}.

\begin{quote}
\mr{इमारतीचे} \mr{दरवाजे} \mr{स्वयंचलित} \mr{आहेत}.
\end{quote}

\begin{quote}
\emph{imārtīce darvāje svayañcalit āhet.}
\end{quote}

\begin{quote}
The building's doors are automatic.
\end{quote}

\begin{cousinweb}
\textbf{\mr{स्वयंचलित}} (\emph{svayañcalit}) is a modern tatsama coinage from Sanskrit \emph{svayam}, \textquotedblleft{}by itself\textquotedblright{}, and \emph{calita}, \textquotedblleft{}moving\textquotedblright{} (Monier-Williams).
\end{cousinweb}

\begin{grammarlens}[title={What you've built}]
One more word for describing what you lived and felt.
\end{grammarlens}

\subsection*{Your turn}
\begin{itemize}
\raggedright
  \item \emph{Say it:} \emph{svayañcalit}
  \item \emph{Say it:} \emph{svayañcalit}, once more
  \item \emph{Say it:} \emph{imārtīce darvāje svayañcalit āhet.} — The building's doors are automatic.
  \item \emph{Recall:} say the Marathi for a device, then the Marathi for a robot, then say \emph{svayañcalit} again
\end{itemize}

\begin{tcolorbox}[breakable,colback=teal!4,colframe=teal!35!black,title={Before you move on}]
What does \mr{स्वयंचलित} mean? (\textquotedblleft{}Automatic\textquotedblright{}.) Say it once more.
\end{tcolorbox}
\section[\texorpdfstring{\mr{फवारणे} (phavārṇe)}{फवारणे (phavārṇe)}]{\mr{फवारणे} (phavārṇe) — to spray}

\label{lesson:MR-C518-phavarne}



\begin{quote}
Before the new one: say the Marathi for a robot, then the Marathi for automatic.
\end{quote}

\subsection*{You'll want to know: \mr{फवारणे}}
\textbf{\mr{फवारणे}} — \emph{phavārṇe} — \textquotedblleft{}to spray\textquotedblright{}.

\begin{quote}
\mr{यंत्राने} \mr{शेतावर} \mr{औषध} \mr{फवारलं}.
\end{quote}

\begin{quote}
\emph{yantrāne śetāvar auṣadh phavārlaṁ.}
\end{quote}

\begin{quote}
The machine sprayed the field with pesticide.
\end{quote}

Past, with -\mr{ने} (he …): \mr{त्याने} \mr{औषध} \mr{फवारलं}, \emph{tyāne auṣadh phavārlaṁ}.

The noun is \mr{फवारा}, a spray or a fountain.

\begin{cousinweb}
\textbf{\mr{फवारणे}} (\emph{phavārṇe}) is built on \mr{फवारा}, \textquotedblleft{}a spray, a fountain\textquotedblright{}, a loan through Persian usually traced to Arabic \emph{fawwāra}, \textquotedblleft{}fountain\textquotedblright{}.
\end{cousinweb}

\begin{grammarlens}[title={What you've built}]
One more word for describing what you lived and felt.
\end{grammarlens}

\subsection*{Your turn}
\begin{itemize}
\raggedright
  \item \emph{Say it:} \emph{phavārṇe}
  \item \emph{Say it:} \emph{phavārṇe}, once more
  \item \emph{Say it:} \emph{yantrāne śetāvar auṣadh phavārlaṁ.} — The machine sprayed the field with pesticide.
  \item \emph{Recall:} say the Marathi for a robot, then the Marathi for automatic, then say \emph{phavārṇe} again
\end{itemize}

\begin{tcolorbox}[breakable,colback=teal!4,colframe=teal!35!black,title={Before you move on}]
What does \mr{फवारणे} mean? (\textquotedblleft{}To spray\textquotedblright{}.) Say it once more.
\end{tcolorbox}
